\chapter{Complexity — Mathlib Bridge}
%%% generated by the blueprint pipeline from TCSlib.Complexity.TuringMachine.MathlibBridge
%%%%%%%%%%%%%%%%%%%%%%%%%%%%%%%%%%%%%%%%%%%%%%%%%%%%%%%%%%%%%%%%%%%%%%%%%%%%%%%
\section{Overview}

This module is the Mathlib-facing layer behind the existence of an effective
machine-code scheme.  It proves the erased machine-code parser and the associated
prefix and canonizing string operations primitive recursive, and converts primitive
recursiveness into an actual machine: Mathlib's verified compiler translates a
suitable partial-recursive program into a four-stack machine, a private
four-work-tape bridging machine simulates that stack machine step by step inside
the project's multi-tape model, and the alphabet-reduction theorem lands the result
in a binary machine.  A finite maximum of halting times at each input length
supplies the time bound; no polynomial bound is claimed.

%%%%%%%%%%%%%%%%%%%%%%%%%%%%%%%%%%%%%%%%%%%%%%%%%%%%%%%%%%%%%%%%%%%%%%%%%%%%%%%
\section{Declarations}

\begin{lemma}[Primitive recursiveness of the unary reader]
\lean{Turing.codePrimUnary}
\label{Turing.codePrimUnary}
\leanok
\uses{Turing.codeReadUnary}
\difficulty{4}
The unary reader on Boolean strings --- the partial function that consumes a maximal
leading block of true bits together with the terminating false bit and returns the
pair consisting of the block's length and the unconsumed suffix, undefined when the
string contains no false bit --- is primitive recursive.
\end{lemma}

\begin{lemma}[Primitive recursiveness of the bit-append operation]
\lean{Turing.codePrimBit}
\label{Turing.codePrimBit}
\leanok
\uses{Turing.codePrimSkipFin}
\difficulty{3}
The two-argument operation taking a Boolean bit $b$ and a natural number $n$ to
$2n + b$, i.e.\ appending $b$ as a new least-significant binary digit, is primitive
recursive.
\end{lemma}

\begin{lemma}[Primitive recursiveness of binary evaluation]
\lean{Turing.codePrimBitsNat}
\label{Turing.codePrimBitsNat}
\leanok
\uses{Turing.codeBitsNat, Turing.codePrimBit}
\difficulty{2}
The function evaluating a Boolean string as a natural number written with the
least-significant bit first (the fold of the bit-append operation
$(b, n) \mapsto 2n + b$ over the string, starting from $0$) is primitive recursive.
\end{lemma}

\begin{lemma}[Primitive recursiveness of the aligned pair parser]
\lean{Turing.codePrimPair}
\label{Turing.codePrimPair}
\leanok
\uses{Turing.pairDecode}
\difficulty{6}
The aligned pair parser is primitive recursive.  This parser reads a Boolean string
in aligned two-bit blocks: the blocks $00$ and $11$ contribute the data bits false
and true, the first aligned block $01$ acts as separator, and any other block (or a
missing separator) is rejected; on success it returns the decoded first component
together with the suffix following the separator.
\end{lemma}

\begin{lemma}[Primitive recursiveness of binary expansion]
\lean{Turing.codePrimBits}
\label{Turing.codePrimBits}
\leanok
\difficulty{6}
The function assigning to each natural number the list of its binary digits,
least-significant bit first (the number $0$ having the empty digit list), is
primitive recursive.
\end{lemma}

\begin{lemma}[Primitive recursiveness of the two-bit skipping reader]
\lean{Turing.codePrimSkipPair}
\label{Turing.codePrimSkipPair}
\leanok
\uses{Turing.codeSkipPair}
\difficulty{3}
For every fixed test $v$ on pairs of Booleans, the skipping reader that consumes the
first two bits $a, b$ of a Boolean string and returns the remaining suffix when
$v(a, b)$ holds --- rejecting otherwise, and rejecting strings with fewer than two
bits --- is primitive recursive.
\end{lemma}

\begin{lemma}[Primitive recursiveness of the bounded-index skipper]
\lean{Turing.codePrimSkipFin}
\label{Turing.codePrimSkipFin}
\leanok
\uses{Turing.codePrimUnary, Turing.codeSkipFin}
\difficulty{3}
The two-argument skipping reader that, given a bound $n$ and a Boolean string,
consumes a unary-encoded index (a block of true bits closed by a false bit) and
returns the unconsumed suffix precisely when the index is smaller than $n$, is
primitive recursive jointly in $n$ and the string.
\end{lemma}

\begin{lemma}[Primitive recursiveness of the state-field skipper]
\lean{Turing.codePrimSkipState}
\label{Turing.codePrimSkipState}
\leanok
\uses{Turing.codePrimSkipFin, Turing.codeSkipFin, Turing.codeSkipState}
\difficulty{3}
The two-argument skipping reader for a halt-or-successor field is primitive
recursive jointly in the bound and the string.  This reader consumes one tag bit;
on a true tag it further consumes a unary index checked to be smaller than the
bound, and on a false tag it returns the rest of the string unchanged.
\end{lemma}

\begin{lemma}[Primitive recursiveness of the record skipper]
\lean{Turing.codePrimSkipAction}
\label{Turing.codePrimSkipAction}
\leanok
\uses{Turing.codePrimSkipPair, Turing.codePrimSkipState, Turing.codeSkipAction}
\difficulty{3}
The two-argument skipping reader for one five-field transition record is primitive
recursive jointly in the state bound $n$ and the string.  The record consists of
four two-bit fields --- the first, third, and fourth checked against a reserved
block, the second unconstrained --- followed by a halt-or-successor field whose
unary index must be smaller than $n + 1$.
\end{lemma}

\begin{lemma}[Iterated skipping as an iterate on optional strings]
\lean{Turing.codeSkipRepeat_iter}
\label{Turing.codeSkipRepeat_iter}
\leanok
\uses{Turing.codeSkipRepeat}
\difficulty{4}
Let $r$ be a partial reader on Boolean strings (returning either a suffix or
failure), let $n \in \bbn$, and let $xs$ be a Boolean string.  Applying $r$
sequentially $n$ times starting from $xs$, threading each returned suffix into the
next application and failing as soon as one application fails, gives the same result
as the $n$-th iterate --- evaluated at $xs$ --- of the total map on optional strings
that feeds a present value to $r$ and propagates failure.
\end{lemma}

\begin{lemma}[Primitive recursiveness of counted iteration]
\lean{Turing.codePrimRepeat}
\label{Turing.codePrimRepeat}
\leanok
\uses{Turing.codeSkipRepeat, Turing.codeSkipRepeat_iter}
\difficulty{3}
Let $A$ be a type with a primitive recursive coding, let $r$ assign to each
$a \in A$ a partial reader on Boolean strings, primitive recursively jointly in $a$
and the string, and let $\mathrm{count} \colon A \to \bbn$ be primitive recursive.
Then the function taking $a$ and a string to the result of applying the reader
$r(a)$ sequentially $\mathrm{count}(a)$ times is primitive recursive jointly in $a$
and the string.
\end{lemma}

\begin{lemma}[Primitive recursiveness of the all-true test]
\lean{Turing.codePrimAll}
\label{Turing.codePrimAll}
\leanok
\difficulty{3}
The Boolean-valued function testing whether every entry of a Boolean string is true
is primitive recursive.
\end{lemma}

\begin{lemma}[Primitive recursiveness of the erased parser]
\lean{Turing.codePrimScan}
\label{Turing.codePrimScan}
\leanok
\uses{Turing.codePrimAll, Turing.codePrimBits, Turing.codePrimBitsNat, Turing.codePrimPair, Turing.codePrimRepeat, Turing.codePrimSkipAction, Turing.codePrimSkipFin, Turing.codeScan}
\difficulty{4}
The erased machine-code parser is primitive recursive.  On a Boolean string this
parser recovers a state count $n$ from the aligned doubled-bit region, rejects
unless those bits are exactly the canonical binary expansion of $n$, rejects when
fewer than $81(n+1)$ bits remain, skips a unary initial state below $n + 1$ and then
$9(n+1)$ five-field transition records, and accepts precisely when the leftover
suffix consists entirely of true bits, returning that suffix.
\end{lemma}

\begin{lemma}[Primitive recursiveness of the drop operation]
\lean{Turing.codePrimDrop}
\label{Turing.codePrimDrop}
\leanok
\difficulty{4}
The two-argument function taking a Boolean string and a natural number $n$ to the
string with its first $n$ entries removed is primitive recursive.
\end{lemma}

\begin{lemma}[Primitive recursiveness of the co-length prefix]
\lean{Turing.codePrimPrefix}
\label{Turing.codePrimPrefix}
\leanok
\uses{Turing.codePrimDrop}
\difficulty{2}
The two-argument function taking a Boolean string $xs$ and a natural number $n$ to
the prefix of $xs$ of length $\abs{xs} - n$ (truncated subtraction, so the empty
string when $n \ge \abs{xs}$) is primitive recursive.
\end{lemma}

\begin{lemma}[Primitive recursiveness of the canonizing function]
\lean{Turing.codePrimCanonical}
\label{Turing.codePrimCanonical}
\leanok
\uses{Turing.CodeTM.serialize, Turing.codeCanonical, Turing.codeFallback, Turing.codePrimPrefix, Turing.codePrimScan}
\difficulty{2}
The canonizing string function is primitive recursive.  This function maps a Boolean
string $xs$ to the prefix consumed by the erased machine-code parser when that
parser accepts (namely the first $\abs{xs} - s$ bits, where $s$ is the length of the
leftover all-true suffix), and to the fixed serialization of the trivial one-state
halting machine when the parser rejects.
\end{lemma}

\begin{definition}[Bridge alphabet]
\lean{Turing.BridgeAlphabet}
\label{Turing.BridgeAlphabet}
\leanok
The alphabet of the bridging machine: the disjoint union of the Boolean input
alphabet with the four-symbol alphabet of Mathlib's compiled stack machines (two bit
symbols and two list-separator symbols).
\end{definition}

\begin{definition}[Enumeration of the four stacks]
\lean{Turing.bridgeIndex}
\label{Turing.bridgeIndex}
\leanok
The enumeration of the four stack names of Mathlib's compiled stack machines ---
main, reverse, auxiliary, and call stack --- as work-tape indices $0, 1, 2, 3$.
\end{definition}

\begin{definition}[Tape representation of a stack]
\lean{Turing.bridgeStack}
\label{Turing.bridgeStack}
\leanok
\uses{Turing.BridgeAlphabet}
The work-tape contents representing a stack: a stack $xs$ over the compiled stack
machine's alphabet occupies the cells $-\abs{xs}, \dots, -1$ of a work tape, with
the top of the stack at cell $-\abs{xs}$, each symbol embedded into the bridge
alphabet, and every other cell blank.
\end{definition}

\begin{lemma}[Reading the top of a represented stack]
\lean{Turing.bridgeStack_read}
\label{Turing.bridgeStack_read}
\leanok
\uses{Turing.bridgeStack}
\difficulty{2}
In the work-tape representation of a stack $xs$, the cell at position $-\abs{xs}$
holds the embedded top symbol of $xs$, and is blank when $xs$ is empty.
\end{lemma}

\begin{lemma}[Representation of the empty stack]
\lean{Turing.bridgeStack_nil}
\label{Turing.bridgeStack_nil}
\leanok
\uses{Turing.bridgeStack}
\difficulty{1}
The work-tape representation of the empty stack is the everywhere-blank tape.
\end{lemma}

\begin{lemma}[Pushing onto a represented stack]
\lean{Turing.bridgeStack_push}
\label{Turing.bridgeStack_push}
\leanok
\uses{Turing.bridgeStack}
\difficulty{4}
Updating the work-tape representation of a stack $xs$ by writing an embedded symbol
$a$ into the cell at position $-\abs{xs} - 1$ yields exactly the representation of
the pushed stack with $a$ prepended to $xs$.
\end{lemma}

\begin{lemma}[Popping a represented stack]
\lean{Turing.bridgeStack_pop}
\label{Turing.bridgeStack_pop}
\leanok
\uses{Turing.bridgeStack, Turing.bridgeStack_push}
\difficulty{4}
Updating the work-tape representation of a nonempty stack, with top symbol $a$ above
a tail $xs$, by blanking its top cell (at position $-(\abs{xs} + 1)$) yields exactly
the representation of the popped stack $xs$.
\end{lemma}

\begin{definition}[Inverse enumeration of the four stacks]
\lean{Turing.bridgeKey}
\label{Turing.bridgeKey}
\leanok
The inverse enumeration of stack names: the map sending the work-tape indices
$0, 1, 2, 3$ to the four stack names main, reverse, auxiliary, and call stack of
Mathlib's compiled stack machines.
\end{definition}

\begin{lemma}[Round trip on stack names]
\lean{Turing.bridgeKey_index}
\label{Turing.bridgeKey_index}
\leanok
\uses{Turing.bridgeIndex, Turing.bridgeKey}
\difficulty{2}
Mapping a stack name to its work-tape index and back returns the original stack
name.
\end{lemma}

\begin{lemma}[Round trip on tape indices]
\lean{Turing.bridgeIndex_key}
\label{Turing.bridgeIndex_key}
\leanok
\uses{Turing.bridgeIndex, Turing.bridgeKey}
\difficulty{3}
Mapping a work-tape index in $\{0, 1, 2, 3\}$ to its stack name and back returns the
original index.
\end{lemma}

\begin{definition}[Control states of the bridging machine]
\lean{Turing.BridgeState}
\label{Turing.BridgeState}
\leanok
The control states of the bridging machine over a type $Q$ of source statements: a
forward input-scanning state; two seeding states that lay down the list terminator
and the sentinel bit; a backward-reading state and a state that pushes one read
input bit; execution and push states carrying a source statement from $Q$ together
with the optional register contents; and an emitting state carrying one optional
pending output bit.
\end{definition}

\begin{definition}[Finite statement support of a compiled program]
\lean{Turing.bridgeSupp}
\label{Turing.bridgeSupp}
\leanok
The finite statement support of a compiled program: for a program $c$ in Mathlib's
partial-recursive language, the finite set of optional stack-machine statements
arising from the labels in the support of $c$'s compilation, together with the halt
marker.
\end{definition}

\begin{definition}[Execution-state index type]
\lean{Turing.BridgeQ}
\label{Turing.BridgeQ}
\leanok
\uses{Turing.bridgeSupp}
The state-index type for the execution phase of the bridging machine of a program
$c$: the elements of the finite statement support of $c$'s compilation, viewed as a
finite type.
\end{definition}

\begin{definition}[Bit symbols]
\lean{Turing.bridgeBit}
\label{Turing.bridgeBit}
\leanok
The embedding of Boolean bits into the compiled stack machine's alphabet, sending
false and true to the two bit symbols.
\end{definition}

\begin{definition}[Optional execution state]
\lean{Turing.bridgeExec}
\label{Turing.bridgeExec}
\leanok
\uses{Turing.BridgeQ, Turing.BridgeState, Turing.bridgeSupp}
The optional execution control state determined by an optional stack-machine
statement $q$ and an optional register value $v$: the execution state carrying $q$
and $v$ when $q$ belongs to the finite statement support of the compiled program,
and no state (immediate halt) otherwise.
\end{definition}

\begin{definition}[Idle action]
\lean{Turing.bridgeIdle}
\label{Turing.bridgeIdle}
\leanok
\uses{Turing.Action, Turing.BridgeAlphabet}
The idle action of the bridging machine: move no head, write nothing, emit no
output, and pass control to a given optional successor state.
\end{definition}

\begin{definition}[One-tape action]
\lean{Turing.bridgeOne}
\label{Turing.bridgeOne}
\leanok
\uses{Turing.Action, Turing.BridgeAlphabet}
A one-tape action of the bridging machine: on a selected work tape, optionally write
a symbol and move that head in a given direction, leaving the input head and all
other work tapes untouched, emitting no output, and passing control to a given
optional successor state.
\end{definition}

\begin{definition}[The bridging machine]
\lean{Turing.bridgeTM}
\label{Turing.bridgeTM}
\leanok
\uses{Turing.BridgeAlphabet, Turing.BridgeQ, Turing.BridgeState, Turing.FinTM, Turing.bridgeBit, Turing.bridgeExec, Turing.bridgeIdle, Turing.bridgeIndex, Turing.bridgeOne}
For a program $c$ in Mathlib's partial-recursive language, the bridging machine of
$c$: a finite Turing machine over the bridge alphabet with four work tapes, one per
stack of the compiled stack machine.  Each stack occupies the negative cells of its
tape ending at $-1$, with the head on the stack top and a blank head at zero for an
empty stack.  The machine scans its input to the right end, seeds the main tape with
the list terminator and a sentinel true bit, reads the input backwards while pushing
each bit onto the main tape, executes the compiled program's statements over the
finite statement support, and finally pops the main stack while emitting the result
bits to the output tape.
\end{definition}

\begin{definition}[Bridge configuration]
\lean{Turing.bridgeCfg}
\label{Turing.bridgeCfg}
\leanok
\uses{Turing.BridgeAlphabet, Turing.BridgeQ, Turing.BridgeState, Turing.Cfg, Turing.bridgeKey, Turing.bridgeStack}
The bridge configuration determined by an optional control state, an input head
position, an assignment $st$ of stacks to the four stack names, and an output
string: each stack $st(k)$ is laid out on the corresponding work tape in the
standard representation, with that tape's head at the top position $-\abs{st(k)}$.
\end{definition}

\begin{lemma}[Work symbols of a bridge configuration]
\lean{Turing.bridgeCfg_read}
\label{Turing.bridgeCfg_read}
\leanok
\uses{Turing.BridgeAlphabet, Turing.BridgeQ, Turing.BridgeState, Turing.Cfg.workTapeSymbols, Turing.bridgeCfg, Turing.bridgeIndex, Turing.bridgeKey_index, Turing.bridgeStack_read}
\difficulty{1}
In a bridge configuration with stack assignment $st$, the symbol read on the work
tape assigned to a stack name $k$ is the embedded top of $st(k)$, or blank when
$st(k)$ is empty.
\end{lemma}

\begin{definition}[Reachability of configurations]
\lean{Turing.bridgeReach}
\label{Turing.bridgeReach}
\leanok
\uses{Turing.Cfg, Turing.MultiTapeTM, Turing.MultiTapeTM.runFrom}
Reachability between configurations of a multi-tape machine on a fixed input: a
configuration $b$ is reachable from $a$ when running the machine from $a$ for some
finite number of steps yields exactly $b$.
\end{definition}

\begin{lemma}[Reflexivity of reachability]
\lean{Turing.bridgeReach_refl}
\label{Turing.bridgeReach_refl}
\leanok
\uses{Turing.Cfg, Turing.MultiTapeTM, Turing.bridgeReach}
\difficulty{1}
Every configuration of a multi-tape machine is reachable from itself, in zero steps.
\end{lemma}

\begin{lemma}[One step is reachable]
\lean{Turing.bridgeReach_step}
\label{Turing.bridgeReach_step}
\leanok
\uses{Turing.Cfg, Turing.MultiTapeTM, Turing.MultiTapeTM.step, Turing.bridgeReach}
\difficulty{1}
The one-step successor of any configuration of a multi-tape machine is reachable
from it.
\end{lemma}

\begin{lemma}[Transitivity of reachability]
\lean{Turing.bridgeReach_trans}
\label{Turing.bridgeReach_trans}
\leanok
\uses{Turing.Cfg, Turing.MultiTapeTM, Turing.MultiTapeTM.runFrom_add, Turing.bridgeReach}
\difficulty{2}
Reachability between configurations of a multi-tape machine on a fixed input is
transitive: if $b$ is reachable from $a$ and $d$ is reachable from $b$, then $d$ is
reachable from $a$.
\end{lemma}

\begin{lemma}[Idle action on a bridge configuration]
\lean{Turing.bridgeCfg_idle}
\label{Turing.bridgeCfg_idle}
\leanok
\uses{Turing.Action.apply, Turing.BridgeAlphabet, Turing.BridgeQ, Turing.BridgeState, Turing.bridgeCfg, Turing.bridgeIdle}
\difficulty{2}
Applying the idle action with a successor state $q'$ to a bridge configuration
changes only the control state to $q'$, leaving the input head, all stacks, and the
output unchanged.
\end{lemma}

\begin{lemma}[Pop action on a bridge configuration]
\lean{Turing.bridgeCfg_pop}
\label{Turing.bridgeCfg_pop}
\leanok
\uses{Turing.Action.apply, Turing.Action.apply_workTapes, Turing.BridgeAlphabet, Turing.BridgeQ, Turing.BridgeState, Turing.bridgeCfg, Turing.bridgeIndex, Turing.bridgeIndex_key, Turing.bridgeKey, Turing.bridgeKey_index, Turing.bridgeOne, Turing.bridgeStack_pop, Turing.moveInputPos_zero}
\difficulty{5}
Suppose that in a bridge configuration the stack assigned to a name $k$ is nonempty,
with top symbol $a$ above a tail $xs$.  Then the one-tape action on the tape of $k$
that blanks the current cell and moves that head one cell to the right, passing to a
control state $q'$, transforms the configuration into the bridge configuration with
control state $q'$ whose assignment replaces the stack at $k$ by $xs$, the input
position and output unchanged.
\end{lemma}

\begin{lemma}[Push as two actions on a bridge configuration]
\lean{Turing.bridgeCfg_push}
\label{Turing.bridgeCfg_push}
\leanok
\uses{Turing.Action.apply, Turing.Action.apply_workTapes, Turing.BridgeAlphabet, Turing.BridgeQ, Turing.BridgeState, Turing.bridgeCfg, Turing.bridgeIndex, Turing.bridgeIndex_key, Turing.bridgeKey, Turing.bridgeKey_index, Turing.bridgeOne, Turing.bridgeStack_push}
\difficulty{5}
Pushing a symbol $a$ onto the stack named $k$ is realized by two consecutive
one-tape actions on the tape of $k$: first move that head one cell to the left
(writing nothing), then write the embedded symbol $a$ without moving.  Applied to a
bridge configuration, this yields the bridge configuration whose assignment replaces
the stack at $k$ by $a$ prepended to it, with the specified final control state and
the input position and output unchanged.
\end{lemma}

\begin{lemma}[Pop of a possibly empty stack]
\lean{Turing.bridgeCfg_pop_any}
\label{Turing.bridgeCfg_pop_any}
\leanok
\uses{Turing.Action.apply, Turing.BridgeAlphabet, Turing.BridgeQ, Turing.BridgeState, Turing.bridgeCfg, Turing.bridgeCfg_pop, Turing.bridgeIndex, Turing.bridgeKey_index, Turing.bridgeOne, Turing.bridgeStack_nil, Turing.moveInputPos_zero}
\difficulty{4}
The one-tape action on the tape of a stack name $k$ that blanks the current cell and
moves the head right precisely when the stack at $k$ is nonempty transforms a bridge
configuration into the bridge configuration whose stack at $k$ is replaced by its
tail (an empty stack remaining empty), with the input position and output unchanged.
\end{lemma}

\begin{lemma}[Execution state of a supported statement]
\lean{Turing.bridgeExec_mem}
\label{Turing.bridgeExec_mem}
\leanok
\uses{Turing.bridgeExec, Turing.bridgeSupp}
\difficulty{1}
If an optional statement $q$ belongs to the finite statement support of the compiled
program, then the optional execution control state determined by $q$ and a register
value $v$ is precisely the execution state carrying $q$ and $v$.
\end{lemma}

\begin{lemma}[Unfolding one step of the bridging machine]
\lean{Turing.bridge_step}
\label{Turing.bridge_step}
\leanok
\uses{Turing.Action.apply, Turing.BridgeAlphabet, Turing.BridgeQ, Turing.BridgeState, Turing.Cfg.inputSymbol, Turing.Cfg.workTapeSymbols, Turing.MultiTapeTM.step, Turing.bridgeCfg, Turing.bridgeTM}
\difficulty{1}
One step of the bridging machine from a bridge configuration in a live control state
is computed by applying the action that its transition table prescribes for that
state, the symbol under the input head, and the symbols under the four work-tape
heads.
\end{lemma}

\begin{lemma}[Support closure under substatements]
\lean{Turing.bridge_sub}
\label{Turing.bridge_sub}
\leanok
\uses{Turing.bridgeSupp}
\difficulty{1}
The finite statement support of a compiled program is closed under immediate
substatements: if a statement belongs to the support and another statement occurs as
an immediate constituent of it, the latter also belongs to the support.
\end{lemma}

\begin{lemma}[The halt marker is supported]
\lean{Turing.bridge_none}
\label{Turing.bridge_none}
\leanok
\uses{Turing.bridgeSupp}
\difficulty{1}
The halt marker belongs to the finite statement support of every compiled program.
\end{lemma}

\begin{lemma}[Simulation of one source step]
\lean{Turing.bridge_statement}
\label{Turing.bridge_statement}
\leanok
\uses{Turing.Action.apply, Turing.BridgeAlphabet, Turing.MultiTapeTM.runFrom_succ_eq_step', Turing.MultiTapeTM.runFrom_zero, Turing.MultiTapeTM.step, Turing.bridgeCfg, Turing.bridgeCfg_idle, Turing.bridgeCfg_pop_any, Turing.bridgeCfg_push, Turing.bridgeCfg_read, Turing.bridgeExec, Turing.bridgeExec_mem, Turing.bridgeIndex, Turing.bridgeOne, Turing.bridgeReach, Turing.bridgeReach_trans, Turing.bridgeSupp, Turing.bridgeTM, Turing.bridge_step, Turing.bridge_sub}
\difficulty{6}
Let $c$ be a program in Mathlib's partial-recursive language, $q$ a stack-machine
statement in the finite statement support of $c$'s compilation, $v$ an optional
register value, and $st$ a stack assignment.  Executing $q$ on $v$ and $st$ in the
stack-machine semantics produces a next optional label, register value, and stacks.
From the bridge configuration in the execution state for $q$ and $v$ over $st$, the
bridging machine of $c$ reaches, in finitely many steps, the bridge configuration in
the optional execution state determined by the translation of that next label and by
that register value, over those resulting stacks, with the input position and output
unchanged.
\end{lemma}

\begin{lemma}[Translated labels are supported]
\lean{Turing.bridge_label}
\label{Turing.bridge_label}
\leanok
\uses{Turing.bridgeSupp, Turing.bridge_none}
\difficulty{3}
For every optional label drawn from the label support of the compiled program with
the halt label adjoined, the statement obtained by translating that label belongs to
the finite statement support of the program.
\end{lemma}

\begin{lemma}[Simulation of compiled runs]
\lean{Turing.bridge_simulate}
\label{Turing.bridge_simulate}
\leanok
\uses{Turing.BridgeAlphabet, Turing.bridgeCfg, Turing.bridgeExec, Turing.bridgeNumber, Turing.bridgeReach, Turing.bridgeReach_refl, Turing.bridgeReach_trans, Turing.bridgeTM, Turing.bridge_label, Turing.bridge_statement}
\difficulty{5}
Let $c$ be a program in Mathlib's partial-recursive language, and let $a, b$ be
configurations of the compiled stack machine such that $b$ is reachable from $a$
and the label of $a$ lies in the label support of $c$'s compilation with the halt
label adjoined.  Then the label of $b$ also lies in that support, and, for every
fixed input head position and output string, the bridging machine of $c$ reaches,
from the bridge configuration associated with $a$ (the optional execution state for
$a$'s translated label and register value, over $a$'s stacks), the bridge
configuration associated with $b$.
\end{lemma}

\begin{definition}[Sentinel numeric code]
\lean{Turing.bridgeNumber}
\label{Turing.bridgeNumber}
\leanok
The sentinel numeric code of a Boolean string: the natural number whose binary
expansion, least-significant bit first, consists of the string's bits followed by
one final true bit.  It is obtained by folding the bit-append operation
$(b, n) \mapsto 2n + b$ over the string, starting from $1$.
\end{definition}

\begin{definition}[Stack word of a Boolean string]
\lean{Turing.bridgeWord}
\label{Turing.bridgeWord}
\leanok
\uses{Turing.bridgeBit}
The stack word of a Boolean string: its bits embedded as bit symbols of the compiled
stack machine's alphabet, followed by a sentinel true-bit symbol and the
list-terminator symbol.
\end{definition}

\begin{lemma}[Positivity of the sentinel code]
\lean{Turing.bridgeNumber_pos}
\label{Turing.bridgeNumber_pos}
\leanok
\uses{Turing.bridgeNumber}
\difficulty{3}
The sentinel numeric code of every Boolean string is positive.
\end{lemma}

\begin{lemma}[Stack encoding of the sentinel code]
\lean{Turing.bridgeWord_number}
\label{Turing.bridgeWord_number}
\leanok
\uses{Turing.bridgeBit, Turing.bridgeNumber, Turing.bridgeNumber_pos, Turing.bridgeWord}
\difficulty{5}
Mathlib's stack encoding of the singleton list containing the sentinel numeric code
of a Boolean string $xs$ equals the stack word of $xs$: the bits of $xs$ as bit
symbols, then the sentinel true-bit symbol, then the list terminator.
\end{lemma}

\begin{lemma}[Push combined with an input move]
\lean{Turing.bridgeCfg_push_input}
\label{Turing.bridgeCfg_push_input}
\leanok
\uses{Turing.Action.apply, Turing.BridgeAlphabet, Turing.BridgeQ, Turing.BridgeState, Turing.bridgeCfg, Turing.bridgeCfg_push, Turing.bridgeIndex, Turing.bridgeOne, Turing.bridgeStore, Turing.moveInputPos}
\difficulty{3}
Performing the two-action push of a symbol $a$ onto the stack named $k$, where the
second action additionally moves the input head in a direction $d$, transforms a
bridge configuration into the bridge configuration with $a$ prepended to the stack
at $k$ and the input head moved by $d$, the output unchanged.
\end{lemma}

\begin{definition}[Main-stack storage]
\lean{Turing.bridgeStore}
\label{Turing.bridgeStore}
\leanok
The stack assignment holding a given word on the main stack and leaving the reverse,
auxiliary, and call stacks empty.
\end{definition}

\begin{lemma}[Pushing onto the main-stack storage]
\lean{Turing.bridgeStore_push}
\label{Turing.bridgeStore_push}
\leanok
\uses{Turing.bridgeStore}
\difficulty{2}
Prepending a symbol to the main stack of the assignment that stores a word on the
main stack (other stacks empty) yields the assignment storing the extended word.
\end{lemma}

\begin{lemma}[Backward input loading]
\lean{Turing.bridge_back}
\label{Turing.bridge_back}
\leanok
\uses{Turing.Action.apply, Turing.BridgeAlphabet, Turing.BridgeQ, Turing.BridgeState, Turing.Cfg.inputSymbol, Turing.MultiTapeTM.runFrom_succ_eq_step', Turing.MultiTapeTM.runFrom_zero, Turing.bridgeBit, Turing.bridgeCfg, Turing.bridgeCfg_idle, Turing.bridgeCfg_push_input, Turing.bridgeExec, Turing.bridgeIndex, Turing.bridgeOne, Turing.bridgeReach, Turing.bridgeReach_trans, Turing.bridgeStore, Turing.bridgeStore_push, Turing.bridgeTM, Turing.bridgeWord, Turing.bridge_step, Turing.inputSymbolInner, Turing.moveInputPos_neg_of_ne_left}
\difficulty{6}
Let $c$ be a program in Mathlib's partial-recursive language and $x$ a Boolean input
string.  For every position $j \le \abs{x}$: from the bridge configuration of the
bridging machine of $c$ in the backward-reading state, with input head at position
$j$ and the main stack holding the stack word of the suffix of $x$ from position $j$
onward (other stacks empty, empty output), the machine reaches the configuration
with the input head at the left endmarker whose main stack holds the stack word of
all of $x$ and whose control is the execution state at the translated initial
statement of $c$'s compilation, with empty register.
\end{lemma}

\begin{lemma}[Empty main-stack storage]
\lean{Turing.bridgeStore_nil}
\label{Turing.bridgeStore_nil}
\leanok
\uses{Turing.bridgeStore}
\difficulty{2}
The stack assignment storing the empty word on the main stack has all four stacks
empty.
\end{lemma}

\begin{lemma}[Evaluating the main-stack storage]
\lean{Turing.bridgeStore_at}
\label{Turing.bridgeStore_at}
\leanok
\uses{Turing.bridgeIndex_key, Turing.bridgeKey, Turing.bridgeStore}
\difficulty{3}
Evaluating the assignment that stores a word $xs$ on the main stack at the stack
name enumerated by an index $i \in \{0, 1, 2, 3\}$ yields $xs$ when $i = 0$ (the
main stack) and the empty stack otherwise.
\end{lemma}

\begin{lemma}[Forward input scan]
\lean{Turing.bridge_scan}
\label{Turing.bridge_scan}
\leanok
\uses{Turing.Action.apply, Turing.BridgeAlphabet, Turing.Cfg.inputSymbol, Turing.MultiTapeTM.initCfg, Turing.MultiTapeTM.runFrom, Turing.MultiTapeTM.runFrom_succ_eq_step', Turing.bridgeCfg, Turing.bridgeStack_nil, Turing.bridgeStore, Turing.bridgeStore_nil, Turing.bridgeTM, Turing.bridge_step, Turing.inputSymbolInner, Turing.moveInputPos, Turing.moveInputPos_pos_of_ne_right}
\difficulty{5}
Let $c$ be a program in Mathlib's partial-recursive language and $x$ a Boolean input
string.  For every $j \le \abs{x}$, running the bridging machine of $c$ for exactly
$j$ steps from its initial configuration on the embedded input $x$ yields the bridge
configuration in the scanning state with the input head at position $j + 1$
(positions are shifted by one, position $0$ being the left endmarker), all stacks
empty, and empty output.
\end{lemma}

\begin{lemma}[Seeding the main stack]
\lean{Turing.bridge_seed}
\label{Turing.bridge_seed}
\leanok
\uses{Turing.Action.apply, Turing.Cfg.inputSymbol, Turing.MultiTapeTM.runFrom, Turing.MultiTapeTM.runFrom_succ_eq_step', Turing.MultiTapeTM.runFrom_zero, Turing.MultiTapeTM.step, Turing.bridgeCfg, Turing.bridgeOne, Turing.bridgeStack, Turing.bridgeStack_push, Turing.bridgeStore, Turing.bridgeStore_at, Turing.bridgeTM, Turing.bridge_step, Turing.moveInputPos_neg_of_ne_left, Turing.moveInputPos_zero}
\difficulty{5}
From the bridge configuration of the bridging machine in the scanning state at the
right endmarker of the embedded Boolean input $x$ (all stacks empty, empty output),
exactly three steps lead to the backward-reading state with the input head at
position $\abs{x}$ and the main stack holding just the sentinel true-bit symbol on
top of the list terminator, the other stacks empty and the output empty.
\end{lemma}

\begin{lemma}[Startup of the bridging machine]
\lean{Turing.bridge_start}
\label{Turing.bridge_start}
\leanok
\uses{Turing.MultiTapeTM.initCfg, Turing.bridgeCfg, Turing.bridgeExec, Turing.bridgeReach, Turing.bridgeReach_trans, Turing.bridgeStore, Turing.bridgeTM, Turing.bridgeWord, Turing.bridge_back, Turing.bridge_scan, Turing.bridge_seed}
\difficulty{3}
From its initial configuration on the embedded Boolean input $x$, the bridging
machine of a partial-recursive program $c$ reaches the bridge configuration at input
position $0$ whose main stack holds the full stack word of $x$ (other stacks empty,
empty output), in the execution state at the translated initial statement of $c$'s
compilation, with empty register.
\end{lemma}

\begin{lemma}[Pop combined with an output emission]
\lean{Turing.bridgeCfg_pop_emit}
\label{Turing.bridgeCfg_pop_emit}
\leanok
\uses{Turing.Action.apply, Turing.BridgeAlphabet, Turing.BridgeQ, Turing.BridgeState, Turing.bridgeCfg, Turing.bridgeCfg_pop, Turing.bridgeIndex, Turing.bridgeOne, Turing.moveInputPos_zero}
\difficulty{3}
Suppose that in a bridge configuration the stack at a name $k$ is nonempty, with top
symbol $a$ above a tail $xs$.  The one-tape pop action on the tape of $k$ (blank the
current cell, move right), extended with the emission of an optional output symbol
$e$, produces the bridge configuration whose stack at $k$ is $xs$ and whose output
is extended by $e$, the input position unchanged.
\end{lemma}

\begin{lemma}[One emission step]
\lean{Turing.bridge_emit_step}
\label{Turing.bridge_emit_step}
\leanok
\uses{Turing.BridgeAlphabet, Turing.Cfg.workTapeSymbols, Turing.MultiTapeTM.step, Turing.bridgeBit, Turing.bridgeCfg, Turing.bridgeCfg_pop_emit, Turing.bridgeCfg_read, Turing.bridgeTM, Turing.bridge_step}
\difficulty{4}
Suppose the bridging machine is in the emitting state with an optional pending bit,
and the main stack begins with the bit symbol embedding a Boolean $b$.  Then a
single step pops that symbol, appends the pending bit (if present, embedded into the
bridge alphabet) to the output, and continues in the emitting state with pending bit
$b$; the input position and the other stacks are unchanged.
\end{lemma}

\begin{lemma}[Emission of the result]
\lean{Turing.bridge_emit}
\label{Turing.bridge_emit}
\leanok
\uses{Turing.BridgeAlphabet, Turing.Cfg.workTapeSymbols, Turing.MultiTapeTM.runFrom_succ_eq_step', Turing.MultiTapeTM.runFrom_zero, Turing.bridgeBit, Turing.bridgeCfg, Turing.bridgeCfg_idle, Turing.bridgeCfg_read, Turing.bridgeReach, Turing.bridgeReach_trans, Turing.bridgeTM, Turing.bridgeWord, Turing.bridge_emit_step, Turing.bridge_step}
\difficulty{5}
Suppose that in a bridge configuration of the bridging machine the main stack holds
the stack word of a Boolean string $xs$ and the control is the emitting state with
an optional pending bit.  Then the machine reaches the halted configuration whose
main stack retains only the list terminator (the other stacks unchanged) and whose
output is extended by the pending bit, if present, followed by the bits of $xs$, all
embedded into the bridge alphabet; the trailing sentinel bit of the stack word is
consumed without being emitted.
\end{lemma}

\begin{lemma}[The bridging machine computes the compiled function]
\lean{Turing.bridge_compiles}
\label{Turing.bridge_compiles}
\leanok
\uses{Turing.FinTM.ComputesInTime, Turing.FinTM.computesInTime_iff, Turing.MultiTapeTM.initCfg, Turing.MultiTapeTM.runFrom, Turing.MultiTapeTM.runFrom_succ_eq_step', Turing.MultiTapeTM.runFrom_zero, Turing.bridgeCfg, Turing.bridgeCfg_idle, Turing.bridgeExec_mem, Turing.bridgeNumber, Turing.bridgeReach, Turing.bridgeReach_trans, Turing.bridgeStore, Turing.bridgeTM, Turing.bridgeWord, Turing.bridgeWord_number, Turing.bridge_emit, Turing.bridge_none, Turing.bridge_simulate, Turing.bridge_start, Turing.bridge_step}
\difficulty{5}
Let $c$ be a program in Mathlib's partial-recursive language and $f$ a total
function on Boolean strings such that, for every string $x$, evaluating $c$ on the
singleton list containing the sentinel numeric code of $x$ yields exactly the
singleton list containing the sentinel numeric code of $f(x)$.  Then for every $x$
there exists a time $t$ such that the bridging machine of $c$, run on the embedded
input $x$, halts within $t$ steps with the embedded string $f(x)$ on its output
tape.
\end{lemma}

\begin{lemma}[A binary machine for the compiled function]
\lean{Turing.bridge_binary}
\label{Turing.bridge_binary}
\leanok
\uses{Turing.FinTM, Turing.FinTM.ComputesFunInTime, Turing.FinTM.ComputesFunInTimeVia, Turing.FinTM.ComputesInTime.mono, Turing.FinTM.alphabet_reduction, Turing.bridgeNumber, Turing.bridgeTM, Turing.bridge_compiles}
\difficulty{4}
Let $c$ be a program in Mathlib's partial-recursive language and $f$ a function on
Boolean strings such that evaluating $c$ on the singleton list holding the sentinel
numeric code of each string $x$ yields exactly the singleton list holding the code
of $f(x)$.  Then there exist a finite Turing machine over the two-letter Boolean
alphabet and a time bound $T \colon \bbn \to \bbn$ such that the machine computes
$f$, halting within $T(\abs{x})$ steps on every input $x$.
\end{lemma}

\begin{lemma}[Binary digits of the sentinel code]
\lean{Turing.bridgeNumber_bits}
\label{Turing.bridgeNumber_bits}
\leanok
\uses{Turing.bridgeNumber, Turing.bridgeNumber_pos}
\difficulty{3}
The binary expansion, least-significant bit first, of the sentinel numeric code of a
Boolean string $xs$ is exactly $xs$ followed by a single true bit.
\end{lemma}

\begin{definition}[Decoding of sentinel codes]
\lean{Turing.bridgeUnnumber}
\label{Turing.bridgeUnnumber}
\leanok
The decoding of sentinel numeric codes: the Boolean string obtained from the binary
expansion of a natural number (least-significant bit first) by removing its final,
most significant digit.
\end{definition}

\begin{lemma}[Decoding inverts the sentinel code]
\lean{Turing.bridgeUnnumber_number}
\label{Turing.bridgeUnnumber_number}
\leanok
\uses{Turing.bridgeNumber, Turing.bridgeNumber_bits, Turing.bridgeUnnumber}
\difficulty{1}
Applying the sentinel decoding to the sentinel numeric code of any Boolean string
returns that string.
\end{lemma}

\begin{lemma}[Primitive recursiveness of the sentinel code]
\lean{Turing.bridgePrimNumber}
\label{Turing.bridgePrimNumber}
\leanok
\uses{Turing.bridgeNumber, Turing.codePrimBit}
\difficulty{2}
The sentinel numeric coding of Boolean strings --- the number whose binary digits
are the string's bits followed by one true bit --- is primitive recursive.
\end{lemma}

\begin{lemma}[Primitive recursiveness of the sentinel decoding]
\lean{Turing.bridgePrimUnnumber}
\label{Turing.bridgePrimUnnumber}
\leanok
\uses{Turing.bridgeUnnumber, Turing.codePrimBits}
\difficulty{2}
The decoding of sentinel numeric codes --- the binary expansion with its most
significant digit removed --- is primitive recursive.
\end{lemma}

\begin{lemma}[Machines for primitive recursive string functions]
\lean{Turing.codePrim_machine}
\label{Turing.codePrim_machine}
\leanok
\uses{Turing.FinTM, Turing.FinTM.ComputesFunInTime, Turing.bridgeNumber, Turing.bridgePrimNumber, Turing.bridgePrimUnnumber, Turing.bridgeUnnumber_number, Turing.bridge_binary}
\difficulty{4}
Every primitive recursive function $f$ on Boolean strings is computed by an actual
machine: there exist a finite Turing machine over the two-letter Boolean alphabet
and a time bound $T \colon \bbn \to \bbn$ such that on every input $x$ the machine
halts within $T(\abs{x})$ steps with $f(x)$ on its output tape.  No growth or
effectivity condition is imposed on $T$.
\end{lemma}

\begin{lemma}[A machine for the canonical serialization]
\lean{Turing.codeCanonical_machine}
\label{Turing.codeCanonical_machine}
\leanok
\uses{Turing.CodeTM.serialize, Turing.FinTM, Turing.FinTM.ComputesFunInTime, Turing.codeCanonical, Turing.codeCanonical_eq, Turing.codeDecode, Turing.codePrimCanonical, Turing.codePrim_machine}
\difficulty{2}
There exist a finite Turing machine over the two-letter Boolean alphabet and a time
bound $T \colon \bbn \to \bbn$ such that on every Boolean string the machine halts
within $T$ of the input's length, outputting the fixed serialization of the
one-work-tape binary machine that the string decodes to (malformed strings decoding
to the trivial one-state halting machine).
\end{lemma}

\begin{theorem}[Existence of an effective machine code]
\lean{Turing.exists_effectiveMachineCode}
\label{Turing.exists_effectiveMachineCode}
\leanok
\uses{Turing.CodeTM.serialize, Turing.EffectiveMachineCode, Turing.codeCanonical_machine, Turing.codeDecode, Turing.codeDecode_serialize_pad}
\difficulty{1}
There exists an effective machine-code scheme: an encoding of one-work-tape binary
machines in code normal form as Boolean strings, and a total decoding, such that a
machine's code followed by any number of trailing true padding bits decodes back to
that machine; together with an actual finite binary Turing machine --- the canonizer
--- and a time bound under which that machine computes, on every Boolean string, the
fixed serialization of the machine the string decodes to.
\end{theorem}
